\chapter{Введение}

\section{Цель}  

Произвести настройку протокола DHCP на маршрутизаторе.
Рассмотреть процесс получения IP – адреса клиентскими устройствами.
\section{Задачи}

\begin{enumerate}
\item Интерфейс Gi1/0/1 маршрутизатора R1 соединить с интерфейсом Gi1/0/1
коммутатора v1ES.
\item Компьютер VPC1 подключить к интерфейсу Gi1/0/10 коммутатора.
\item Компьютер VPC2 подключить к интерфейсу Gi1/0/11 коммутатора.
\item Компьютер VPC3 подключить к интерфейсу Gi1/0/12 коммутатора.
\item Настроить имя маршрутизатору и назначить IP - адрес на интерфейс gi1/0/1
согласно таблице 6.1.1.
\item Настроить DHCP сервер на маршрутизаторе:
\item Создать пул с именем DHCP-10
\item Выдавать из пула диапазон адресов 192.168.10.100 - 192.168.10.200
\item Выдавать адрес шлюза по умолчанию 192.168.10.1
\item Настроить время аренды адреса на 1 час
\item Настроить получение IP - адресов на компьютеры VPC1 – VPC3 по протоколу
DHCP
\end{enumerate}

\begin{figure}[H] 
    \centering
    \includegraphics[width=0.75\textwidth]{top.png}
    \caption{Пример собранной сети}
\end{figure}

\begin{figure}[H] 
    \centering
    \includegraphics[width=0.75\textwidth]{table.png}
    \caption{Таблица с настройками}
\end{figure}

\endinput